\documentclass[11pt,a4paper]{article}

\usepackage[utf8]{inputenc}
\usepackage[T1]{fontenc}
\usepackage[margin=2.5cm]{geometry}
\usepackage{amsmath,amssymb,amsthm}
\usepackage{booktabs}
\usepackage{array}
\usepackage{enumitem}
\usepackage{xcolor}
\usepackage{hyperref}

% ---------------------------------------------------------------------------
% Notation macros (same as docs/batching_formulation.tex)
% ---------------------------------------------------------------------------
\newcommand{\Jset}{\mathcal{J}}
\newcommand{\Oset}{\mathcal{O}}
\newcommand{\Mset}{\mathcal{M}}
\newcommand{\Fset}{\mathcal{F}}
\newcommand{\Bset}{\mathcal{B}}
\newcommand{\Eset}{\mathcal{E}}
\newcommand{\Vset}{\mathcal{V}}
\newcommand{\job}{\operatorname{job}}
\newcommand{\cat}{\,\Vert\,}
\newcommand{\R}{\mathbb{R}}
\newcommand{\Z}{\mathbb{Z}}
\newcommand{\loc}{\ell}
\emergencystretch=3em
\newcommand{\crit}[1]{\par\smallskip\noindent\textbf{Critical note.}\ #1\par\smallskip}

\title{Adding Transport Times to the FJSP with Batching:\\
Formulation and Graph Representations}
\author{Mariana Correia}
\date{\today}

\begin{document}
\maketitle

\noindent This document extends \texttt{docs/batching\_formulation.tex}. It
uses the same notation and refers to its equations as
(F1)--(F2), (B1)--(B7'). Section~\ref{sec:critique} collects the critical
assessment of the design and the open decisions.

\tableofcontents

% ===========================================================================
\section{Transport constraint}
\label{sec:problem}

\subsection{Shop-floor layout}
Every machine $\mu\in\Mset$ has a position $(X_\mu,Y_\mu)\in\Z_{\ge0}^2$ on the
shop floor. Coordinates are measured in \emph{time units}, the same unit as the
processing times $p_{o,\mu}$. There is also one load/unload station (depot),
denoted $\mu=0$, with position $(X_0,Y_0)$. Let $\Mset_0=\Mset\cup\{0\}$.
The transport time between two locations is their Manhattan distance:
\begin{equation}
  \tau_{\mu\nu}=|X_\mu-X_\nu|+|Y_\mu-Y_\nu|,\qquad \mu,\nu\in\Mset_0 .
  \label{eq:tau}
\end{equation}
$\tau$ is a metric: $\tau_{\mu\mu}=0$, $\tau_{\mu\nu}=\tau_{\nu\mu}$ and
$\tau_{\mu\nu}\le\tau_{\mu\kappa}+\tau_{\kappa\nu}$. Because of the triangle
inequality, routing a job through an intermediate location never shortens its
transport. Manhattan (rather than Euclidean) distance models material that moves
along perpendicular aisles. With integer coordinates every $\tau_{\mu\nu}$ is an
integer, which keeps the CP-SAT model exact.

\subsection{Assumptions}
\begin{enumerate}[label=(T\arabic*)]
  \item \label{T:unlimited}\textbf{Unlimited transport capacity.} There are no
    vehicles (AGVs, robots or operators) to schedule. Any number of jobs can be
    transported at the same time. Transport is therefore a
    \emph{machine-dependent time lag} between consecutive operations of a job.
  \item \label{T:immediate}Transport starts as soon as the operation (or its
    batch) finishes. Loading and unloading times are zero or are included in $\tau$.
  \item \label{T:buffer}Buffers are unlimited. A machine becomes free when its
    batch finishes, because the transport does not occupy it. A job that arrives
    before its machine is free waits in the buffer.
  \item \label{T:same}If two consecutive operations of a job run on the same
    machine, no transport is needed ($\tau_{\mu\mu}=0$).
  \item \label{T:depot}Every job starts at the depot. The transport from the
    depot to the machine of its first operation is counted. The return of the
    finished job to the depot is \emph{not} counted in the makespan (see
    Section~\ref{sec:critique}, decision D2).
  \item \label{T:batch}Batch members are transported \emph{individually}. After
    a batch finishes on $\mu$, all of its members are located at $\mu$, and each
    one travels to the machine of its own next operation.
\end{enumerate}

For an operation $o$ let $\mu(o)$ be the machine that processes it. The
transport constraint replaces the routing precedence (B6) by
\begin{equation}
  S_{O_{j,k+1}}\ \ge\ C_{O_{j,k}}+\tau_{\mu(O_{j,k}),\,\mu(O_{j,k+1})},
  \qquad
  S_{O_{j,1}}\ \ge\ \tau_{0,\,\mu(O_{j,1})}.
  \label{eq:transport}
\end{equation}
Consequently, a batch $b$ on machine $\mu$ cannot start before its last member
has arrived:
\begin{equation}
  S_b\ \ge\ \max_{o\in b}\Bigl(C_{\mathrm{pred}(o)}+\tau_{\mu(\mathrm{pred}(o)),\,\mu}\Bigr),
  \label{eq:batcharrival}
\end{equation}
where $\mathrm{pred}(o)$ is the previous operation of $\job(o)$, with $C=0$ and
location $0$ for a first operation.

% ===========================================================================
\section{Mixed-integer formulation}
\label{sec:milp}

The batching model is kept unchanged. Only the machine of each operation and
the transport lags are added.

\paragraph{Additional variables.}
$w_{o,\mu}\in\{0,1\}$: operation $o$ is processed on machine $\mu\in\Mset_o$.

\paragraph{Additional constraints.}
\begin{align}
  & w_{o,\mu}\ \ge\ x_{o,b}+y_{b,\mu}-1
    && \forall o\in\Oset,\ b\in\Bset_{\varphi(o)},\ \mu\in\Mset_{\varphi(o)}
    \tag{T-a}\label{eq:Ta}\\
  & \textstyle\sum_{\mu\in\Mset_o} w_{o,\mu}=1,\qquad w_{o,\mu}=0\ \ \forall\mu\notin\Mset_o
    && \forall o\in\Oset
    \tag{T-b}\label{eq:Tb}\\
  & S_{O_{j,k+1}}\ \ge\ C_{O_{j,k}}+\tau_{\mu\nu}\bigl(w_{O_{j,k},\mu}+w_{O_{j,k+1},\nu}-1\bigr)
    && \forall j,\ k<n_j,\ \mu\neq\nu
    \tag{T-c}\label{eq:Tc}\\
  & S_{O_{j,1}}\ \ge\ \textstyle\sum_{\nu\in\Mset_{O_{j,1}}}\tau_{0\nu}\,w_{O_{j,1},\nu}
    && \forall j\in\Jset
    \tag{T-d}\label{eq:Td}
\end{align}
In \eqref{eq:Tc}, $\mu\in\Mset_{O_{j,k}}$ and $\nu\in\Mset_{O_{j,k+1}}$.
The big constant becomes
$L=\sum_{o}\max_\mu p_{o,\mu}+\sum_j n_j\max_{\mu,\nu\in\Mset_0}\tau_{\mu\nu}$.

\paragraph{Interpretation.}
\begin{itemize}
  \item \eqref{eq:Ta}--\eqref{eq:Tb}: by (B1) and (B4), exactly one batch $b$
    contains $o$ and exactly one machine $\mu^\ast$ has $y_{b,\mu^\ast}=1$.
    \eqref{eq:Ta} forces $w_{o,\mu^\ast}=1$ and \eqref{eq:Tb} forces all other
    $w_{o,\mu}$ to $0$. So $w$ is exactly the machine of $o$, without a
    bilinear term.
  \item \eqref{eq:Tc}: if $O_{j,k}$ runs on $\mu$ and $O_{j,k+1}$ on $\nu$, the
    bracket equals $1$ and the lag $\tau_{\mu\nu}$ is enforced. Otherwise the
    bracket is $\le 0$ and the constraint is weaker than (B6), which is kept.
    No big-$M$ is needed. The number of constraints is
    $\sum_j\sum_{k<n_j}|\Mset_{O_{j,k}}|\,|\Mset_{O_{j,k+1}}|\le|\Oset|\,m^2$.
  \item \eqref{eq:Td}: transport from the depot to the first machine.
  \item \eqref{eq:batcharrival} does not need its own constraint. It follows
    from (B5a), which sets $S_o=S_b$ for every member, together with \eqref{eq:Tc}.
\end{itemize}
With $\tau\equiv0$ (all machines and the depot at one point) the model reduces
exactly to the batching model. This is also the main regression test for the
implementation (Section~\ref{sec:impl}).

\paragraph{CP-SAT version.}
In \texttt{src/batch\_solver.py} each potential batch already has one presence
literal per eligible machine. For every operation $o$ and machine $\mu$, a
literal $z_{o,\mu}$ is created with
$x_{o,b}\wedge y_{b,\mu}\Rightarrow z_{o,\mu}$ and
$\sum_\mu z_{o,\mu}=1$. Then \eqref{eq:Tc} becomes
\begin{center}\small
\texttt{model.add(S[c] >= C[a] + tau[mu][nu])}\\
\texttt{\ \ \ \ .only\_enforce\_if([z[a,mu], z[c,nu]])}
\end{center}
The horizon must include the transport term of $L$.

% ===========================================================================
\section{Consequences for the decision process}
\label{sec:mdp}

\subsection{State and dynamics}
The environment keeps, for every active job $j$, its completion time $C_j$ (the
end of its last scheduled operation, $0$ at the start) and its \emph{location}
$\loc_j\in\Mset_0$ ($\loc_j=0$ at the start). The arrival time of job $j$ at
machine $\mu$ and its earliest start there are
\begin{equation}
  a_{j\mu}=C_j+\tau_{\loc_j\mu},\qquad
  s_{j\mu}=\max\bigl(F_\mu,\ a_{j\mu}\bigr),
  \label{eq:arrival}
\end{equation}
where $F_\mu$ is the free time of $\mu$. Without transport, $a_{j\mu}=C_j$ for
every $\mu$. With transport, \textbf{readiness depends on the machine}. This one
change affects most of the batching machinery.

\paragraph{Batch completion rule.} The selected pair $(\mu,j)$ opens a batch
starting at $S_{b_t}=s_{j\mu}$. A candidate $j'$ now joins only if it has
\emph{arrived at $\mu$} by then:
\begin{equation}
  \varphi(o_{j'}^t)=f,\quad \mu\in\Mset_{o_{j'}^t},\quad
  a_{j'\mu}\le S_{b_t},
\end{equation}
together with the capacity and order-difference conditions, adding candidates in
order of arrival time $a_{j'\mu}$ (previously: ready time). The property
``joining never delays the batch'' is kept.

\paragraph{Update after a decision.} For every member $j$ of the batch:
$C_j\leftarrow S_{b_t}+p_{b_t}$ and $\loc_j\leftarrow\mu$. For every other job,
$s_{j\mu}$ is updated through $F_\mu$ as before.

\paragraph{Reward, mask and teacher.} The reward (makespan increase) is
unchanged. The mask and the warm-start teacher keep their definitions (earliest
start, tie-break earliest completion). Because they are computed from
$s_{j\mu}$ in \eqref{eq:arrival}, they are \emph{transport-aware} in all
representations, including the one without transport information.

\subsection{Batching features under transport}
The batching features of the batching document assume a machine-independent
ready time $r_j$. They must be re-evaluated at a specific machine. Let
\begin{equation}
  \mu^\ast_j=\arg\min_{\mu\in\Mset_{o_j^t}} s_{j\mu}
  \quad(\text{ties: smallest } s_{j\mu}+p_{o_j^t,\mu}),
\end{equation}
be the machine at which job $j$ can start earliest. Then
\begin{multline}
  \pi_j(t)=\min\Bigl(B_{\varphi(o_j^t)}-1,\ \bigl|\{j'\neq j:\ \varphi(o_{j'}^t)=\varphi(o_j^t),\
  |\kappa(o_{j'}^t)-\kappa(o_j^t)|\le\Delta,\\
  \mu^\ast_j\in\Mset_{o_{j'}^t},\
  a_{j'\mu^\ast_j}\le s_{j\mu^\ast_j}\}\bigr|\Bigr).
\end{multline}
This matches the batch completion rule at $\mu^\ast_j$ exactly. The old condition
$\Mset_{o_j^t}\cap\Mset_{o_{j'}^t}\neq\emptyset$ is replaced by the stricter
$\mu^\ast_j\in\Mset_{o_{j'}^t}$. For the family node (Representation A), $t_f$ is
the smallest $s_{j\mu^\ast_j}$ over the current members, $\mu_f$ the machine
that achieves it, and $R_f(t)=\{o_{j'}^t\in\Oset^f:\ \mu_f\in\Mset_{o_{j'}^t},\
a_{j'\mu_f}\le t_f\}$.

Without transport ($\tau\equiv0$) readiness does not depend on the machine, and the
original definitions of the batching document are kept unchanged. This keeps the
batching-only environments, and the models already trained on them, exactly as
they were.

\crit{These features are now \emph{approximations}: they describe the batch at
one machine ($\mu^\ast_j$ or $\mu_f$), while the policy may choose another.
Exact per-machine values belong on per-machine edges. In Representation~A the
\textsf{batch\_exec} edge $(f,\mu)$ has three unused (zero) columns, which could
hold $n_{f\mu}(t)/B_f$ and the earliest batch start at $\mu$. That, however,
changes the batching representation between the experiments with and without
transport. For this reason it is recommended only as a later refinement.}

\subsection{Action-edge features}
The action edge $(\mu,j)$ keeps its five columns
$[p_b,\ \text{rel.\ speed},\ s_{j\mu},\ s_{j\mu}+p_b,\ s_{j\mu}-F_\mu]$, now with the
transport-aware $s_{j\mu}$ of \eqref{eq:arrival} and the batch rule above. A
sixth column is added (Section~\ref{sec:repTF}):
\begin{equation}
  \mathbf{a}_{\mu j}=\bigl[\ p_b,\ \tfrac{p_{o_j^t,\mu}}{\sum_{\mu'}p_{o_j^t,\mu'}},\
  s_{j\mu},\ s_{j\mu}+p_b,\ s_{j\mu}-F_\mu,\ \theta\,\tau_{\loc_j\mu}\ \bigr]\in\R^6,
  \label{eq:actionedge}
\end{equation}
with $\theta=1$ in the feature representation and $\theta=0$ otherwise. Every
attributed relation of the transport environments therefore has dimension $6$:
the other relations are padded with zero columns. The edge width is a property
of the environment (\texttt{EDGE\_DIM}), and the GNN reads it when it is built. The
batching-only environments keep width $5$.
Column 5, $s_{j\mu}-F_\mu$, now also counts machine idle time \emph{caused by
transport}, i.e.\ the job has not yet arrived.

% ===========================================================================
\section{Transport information to represent}
\label{sec:info}

For the comparison to be fair, both representations carry the \emph{same}
information and differ only in how it is encoded, as Representations A and B
did for batching. Four pieces of information are identified:
\begin{description}[style=nextline]
  \item[I1 -- Layout.] Positions of the machines and distances between them (static).
  \item[I2 -- Location.] Where each active job is now, $\loc_j$ (dynamic).
  \item[I3 -- Immediate transport.] The transport each candidate action would
    cause, $\tau_{\loc_j\mu}$ for $\mu\in\Mset_{o_j^t}$ (dynamic, job--machine pair).
  \item[I4 -- Future transport.] The transport that the pending operations will
    require, given their eligible machines (static for pending operations).
\end{description}
For I4, define for two consecutive operations $o=O_{j,k}$ and $o^+=O_{j,k+1}$
\begin{equation}
  \tau^{\min}_{o\to o^+}=\min_{\mu\in\Mset_o,\ \nu\in\Mset_{o^+}}\tau_{\mu\nu},\qquad
  \tau^{\mathrm{avg}}_{o\to o^+}=\frac{1}{|\Mset_o|\,|\Mset_{o^+}|}
  \sum_{\mu\in\Mset_o}\sum_{\nu\in\Mset_{o^+}}\tau_{\mu\nu}.
  \label{eq:tauexp}
\end{equation}
$\tau^{\min}$ is the unavoidable transport (if all machines have distinct
positions, it is positive exactly when the two operations share no eligible
machine). $\tau^{\mathrm{avg}}$ is the transport
expected under a machine choice that ignores the layout. For the first
operation, $\Mset_o$ is replaced by $\{0\}$ (the depot).

\begin{table}[h]
\centering
\caption{Where each piece of information lives in each transport representation.}
\label{tab:info}
\begin{tabular}{l p{4.3cm} p{4.3cm} l}
\toprule
 & T-F: features & T-E: auxiliary edges & T-0 \\
\midrule
I1 layout     & machine $(X_\mu,Y_\mu)$ + $\tau$ on \textsf{listens} edges & \textsf{transport} relation machine$\to$machine & -- \\
I2 location   & job $(X_{\loc_j},Y_{\loc_j})$ & \textsf{at} relation job$\to$machine & -- \\
I3 immediate  & 6th column of the action edge & \textsf{reach} relation job$\to$machine & implicit in $s_{j\mu}$ \\
I4 future     & operation features $\tau^{\mathrm{in}}$ & \textsf{transfer} relation op$\to$op & -- \\
\bottomrule
\end{tabular}
\end{table}

% ===========================================================================
\section{Transport representation T-0: none}
\label{sec:repT0}

The graph is the one of the chosen batching representation. There is no position,
no location and no distance, and the sixth action-edge column is $0$. The
environment, the mask and the teacher still apply the transport
constraint~\eqref{eq:transport}.

\crit{T-0 is \emph{not} blind to transport. The earliest start $s_{j\mu}$ on
the action edges already includes $\tau_{\loc_j\mu}$, because the start time is
a property of the action rather than of the representation. Hiding it (showing
$\max(F_\mu,C_j)$ instead) would make the edges report start times that are
wrong, and the teacher, which ranks by the true $s_{j\mu}$, could no longer be
imitated. That is the problem already fixed for batching. The consequence is
that T-F and T-E are measured against a baseline that already sees the
\emph{current} effect of transport. What they add is I1, I2, I4 and the
\emph{isolated} value of I3 (the start time mixes transport with machine
availability). The expected gain is therefore smaller than for batching, and
the hypothesis should be stated as: ``explicit transport information improves
decisions whose consequences appear in \emph{later} transports''.}

% ===========================================================================
\section{Transport representation T-F: features}
\label{sec:repTF}

No node type and no relation is added. Transport information is stored as
features of the existing nodes and edges.

\paragraph{Machine node.}
\begin{equation}
  \mathbf{x}^{\mathrm{T}}_\mu=\bigl[\ \mathbf{x}_\mu \cat \tilde X_\mu \cat \tilde Y_\mu\ \bigr].
\end{equation}

\paragraph{Job node.} The position of the machine where the job currently is
(the machine of its last processed operation, or the depot):
\begin{equation}
  \mathbf{x}^{\mathrm{T}}_j(t)=\bigl[\ \mathbf{x}_j(t) \cat \tilde X_{\loc_j} \cat \tilde Y_{\loc_j}\ \bigr].
\end{equation}

\paragraph{Operation node.} The transport \emph{into} the operation:
\begin{equation}
  \mathbf{x}^{\mathrm{T}}_o(t)=\bigl[\ \mathbf{x}_o \cat \tau^{\mathrm{in},\min}_o(t) \cat \tau^{\mathrm{in},\mathrm{avg}}_o(t)\ \bigr],
\end{equation}
where
\begin{equation}
  \tau^{\mathrm{in},\cdot}_o(t)=
  \begin{cases}
    \tau^{\cdot}_{\mathrm{pred}(o)\to o} & \text{if } \mathrm{pred}(o) \text{ is not scheduled},\\[2pt]
    \min / \operatorname{avg}_{\nu\in\Mset_o}\tau_{\loc_{\job(o)}\nu} & \text{if } o \text{ is the current operation of its job},
  \end{cases}
\end{equation}
with $\tau^{\cdot}$ from~\eqref{eq:tauexp}. The feature is static for pending
operations and becomes exact (known origin) when the operation becomes current.

\crit{The original idea was an operation feature ``describing what was
processed''. It would carry no information: scheduled operations are
\emph{removed} from the graph (\texttt{\_remove\_operations}), so the network
never sees a processed operation. What was processed, and where, is kept by the
job node through $\loc_j$. The operation feature is therefore redefined to look
ahead, which is where the operation node is useful.}

\paragraph{Machine--machine edges.} The existing fully connected relation
$(\text{machine},\textsf{listens},\text{machine})$ gets an attribute:
\begin{equation}
  \mathbf{a}_{\mu\nu}=\bigl[\ \tau_{\mu\nu},\ 0,\ 0,\ 0,\ 0,\ 0\ \bigr].
\end{equation}

\paragraph{Action edges.} The sixth column of~\eqref{eq:actionedge} holds
$\tau_{\loc_j\mu}$ ($\theta=1$).

\paragraph{Normalisation of the coordinates.}
$\tilde X=(X-X_{\min})/W$ and $\tilde Y=(Y-Y_{\min})/W$, with
$W=\max(X_{\max}-X_{\min},\,Y_{\max}-Y_{\min})$ computed over $\Mset_0$.
The scale is the same for the two axes and for jobs and machines. These columns
must be \emph{excluded} from the per-column min--max normalisation of
\texttt{normalize\_state}.

\crit{\texttt{normalize\_state} rescales every column to $[-1,1]$ separately
for each node type. The job $x$-coordinate would be scaled by the extent of the
\emph{job locations} and the machine $x$-coordinate by the extent of the
\emph{machines}, so the same physical point would get two different values, and
the network could not relate a job to a machine by position. The shared scale
$W$ avoids this, and one scale for both axes keeps the Manhattan geometry.}

\crit{Two limitations of the features should be discussed in the dissertation.
(i) Absolute coordinates are not translation invariant. Only differences matter,
so the coordinates are mainly useful when combined, e.g.\ in the actor, which
concatenates job and machine embeddings. This is why I3 is also given
explicitly. (ii) The distance on \textsf{listens} is a deterministic function of
the two endpoint coordinates and says nothing about any job. It helps a machine
weigh its neighbours, e.g.\ where its jobs will go next, but the quantity the
decision actually needs, $\tau_{\loc_j\mu}$, is a job--machine value. Strictly
speaking, T-F is ``features on existing nodes \emph{and} existing edges''.}

% ===========================================================================
\section{Transport representation T-E: auxiliary edges}
\label{sec:repTE}

No node feature is added. Transport information is encoded by new relations
whose attributes are transport times. The same information is used as in T-F.
\begin{align}
  \Eset_{MM}^{\mathrm{tr}} &= \{(\mu,\nu):\ \mu\neq\nu,\ \exists j,k:\ \mu\in\Mset_{O_{j,k}},\ \nu\in\Mset_{O_{j,k+1}}\},
  \label{eq:Etr}\\
  \Eset_{JM}^{\mathrm{at}}(t) &= \{(j,\loc_j):\ j \text{ active},\ \loc_j\neq 0\},
  \label{eq:Eat}\\
  \Eset_{JM}^{\mathrm{reach}}(t) &= \{(j,\mu):\ j \text{ active},\ \mu\in\Mset_{o_j^t}\},
  \label{eq:Ereach}\\
  \Eset_{OO}^{\mathrm{tf}}(t) &= \{(o^+,o):\ o=O_{j,k},\ o^+=O_{j,k+1},\ o \text{ not scheduled}\}.
  \label{eq:Etf}
\end{align}

\begin{table}[h]
\centering
\caption{New relations of T-E and their attributes (padded with zeros to dimension 6).}
\label{tab:TE}
\begin{tabular}{l l l l}
\toprule
Edge set & Relation & Attribute & Info \\
\midrule
\eqref{eq:Etr}   & (machine, \textsf{transport}, machine) & $[\tau_{\mu\nu}]$ & I1 \\
\eqref{eq:Eat}   & (job, \textsf{at}, machine) + reverse \textsf{hosts} & none & I2 \\
\eqref{eq:Ereach}& (job, \textsf{reach}, machine) & $[\tau_{\loc_j\mu}]$ & I3 \\
\eqref{eq:Etf}   & (operation, \textsf{transfer}, operation) & $[\tau^{\min}_{o\to o^+},\ \tau^{\mathrm{avg}}_{o\to o^+}]$ & I4 \\
\bottomrule
\end{tabular}
\end{table}
\noindent The action edge keeps $\theta=0$.

\paragraph{Interpretation.}
\begin{itemize}
  \item \textsf{transport} (I1) is a \emph{sparse} layout graph. Only machine
    pairs that some routing can actually use are connected, unlike the fully
    connected \textsf{listens}. Both directions are included ($\tau$ is symmetric).
  \item \textsf{at}/\textsf{hosts} (I2) gives the location as structure. A
    machine learns which jobs are waiting there, which matters for batching,
    because members of the same batch leave from the same machine
    (assumption~\ref{T:batch}). A job at the depot has no \textsf{at} edge. Its
    origin is still visible through the \textsf{reach} attributes.
  \item \textsf{reach} (I3) connects the same pairs as the action edges but is a
    separate relation with its own attention. The actor reads edge attributes
    \emph{only} from the action edges, so in T-E the value $\tau_{\loc_j\mu}$
    reaches the decision only through the embeddings. This is the real
    difference from T-F for I3: T-F gives it to the scorer directly, T-E makes
    the GNN use it.
  \item \textsf{transfer} (I4) runs parallel to \textsf{prec} and in the same
    direction (successor $\to$ predecessor). The current operation therefore
    receives the transport of the downstream part of its routing, like the
    transport-weighted conjunctive arcs of a disjunctive graph.
\end{itemize}

\paragraph{Size.}
$|\Eset^{\mathrm{tr}}_{MM}|\le m(m-1)$, $|\Eset^{\mathrm{at}}_{JM}|\le n$,
$|\Eset^{\mathrm{reach}}_{JM}|$ equals the number of action edges ($\le nm$), and
$|\Eset^{\mathrm{tf}}_{OO}|\le|\Oset|-n$. The graph grows linearly in the
instance size.

\paragraph{Receptive field.} With the two GAT layers used, every piece of
information reaches the scored nodes: I2 and I3 in one hop into the job and
machine of the action, I4 in one hop into the current operation and a second hop
into its job, and I1 in one hop between machines.

\crit{The heterogeneous model aggregates relations with a \emph{mean}
(\texttt{to\_hetero(..., aggr='mean')} in \texttt{src/bopo.py}). Adding
\textsf{at}, \textsf{reach} and \textsf{hosts} changes the number of relations
arriving at job and machine nodes, and therefore the weight of every existing
message. Part of any difference between T-E and T-F may come from this dilution
rather than from the transport information. The same confound exists between
batching Representations A and B. It should be named as a threat to validity,
or controlled by one ablation that adds the T-E relations with all attributes
set to zero.}

\crit{\textsf{at} and \textsf{reach} partly overlap: given the edge \textsf{at}
and the \textsf{transport} attributes, $\tau_{\loc_j\mu}$ is determined. They
are kept separate because a 2-layer GNN cannot reliably reconstruct a pairwise
distance from averaged embeddings. If edges have to be saved, \textsf{at} is the
first candidate to remove.}

% ===========================================================================
\section{Experimental grid}
\label{sec:grid}

Batching and transport are represented independently, which gives nine
configurations:
\begin{center}
\begin{tabular}{l ccc}
\toprule
 & T-0 & T-F & T-E \\
\midrule
batching baseline (\texttt{ojmb\_base}) & \texttt{base+t0} & \texttt{base+tf} & \texttt{base+te} \\
Representation A (\texttt{ojmb\_node})  & \texttt{node+t0} & \texttt{node+tf} & \texttt{node+te} \\
Representation B (\texttt{ojmb\_edge})  & \texttt{edge+t0} & \texttt{edge+tf} & \texttt{edge+te} \\
\bottomrule
\end{tabular}
\end{center}
All nine configurations solve \emph{the same problem} (batching + transport) on
the \emph{same instances}, with the same action space, batch rule, mask, teacher
and reward. They differ only in the graph. \texttt{base+t0} is the full
baseline.

\begin{itemize}
  \item \textbf{Main effects.} The transport effect is read along the rows and
    the batching effect along the columns. Because all configurations are
    evaluated on the same instances, the comparisons are \emph{paired}.
    Friedman's test over the nine configurations, followed by pairwise Wilcoxon
    signed-rank tests with a multiple-comparison correction, is recommended.
  \item \textbf{Interaction.} If the best transport representation depends on
    the batching representation, this is itself a result (e.g.\ \textsf{hosts}
    edges may help most together with the family node).
  \item \textbf{Seeds.} At least three training seeds per configuration. With
    nine configurations this is 27 training runs, so the per-run budget has to
    be fixed before the grid is started.
\end{itemize}

\paragraph{Instance generation and calibration.}
Coordinates are drawn uniformly from $\{0,\dots,L\}^2$, with the depot drawn the
same way. For points uniform on $[0,L]^2$, $\mathbb{E}|X_\mu-X_\nu|=L/3$, so the
expected transport time is $\mathbb{E}[\tau]=2L/3$. To make transport a fraction
$\rho$ of the mean processing time $\bar p$,
\begin{equation}
  L=\left\lceil \tfrac{3}{2}\,\rho\,\bar p\right\rceil,
  \qquad\text{e.g. } \bar p\approx 50,\ \rho=0.3\ \Rightarrow\ L=23.
\end{equation}

\crit{If $\rho$ is too small, transport is irrelevant and every configuration
will perform the same, which cannot be told apart from ``the representation does
not help''. If $\rho$ is too large, transport dominates and the batching signal
is lost. Before training, the CP-SAT solutions for a few values of $\rho$ (e.g.\
$0.1,0.3,0.5$) should be used to report (i) the relative increase of the
makespan over the transport-free instance and (ii) the fraction of job time
spent in transport. The value of $\rho$ is then chosen so that transport is
clearly present but not dominant. The same statistics also go in the
dissertation, to show that the problem actually changed.}

% ===========================================================================
\section{Metrics that must change}
\label{sec:metrics}

\paragraph{Feasibility check.} Add the transport constraint to
\texttt{check\_schedule} (batching solver and \texttt{src/metrics}):
$S_{c}\ge C_{a}+\tau_{\mu(a)\mu(c)}$ for consecutive $a,c$, and
$S_{O_{j,1}}\ge\tau_{0\mu(O_{j,1})}$.

\paragraph{Lower bound.} The machine-set bound (b) of \texttt{lower\_bound} is
still valid, because transport does not occupy machines. The job bound (a) must
include transport, otherwise it is loose and the relative error to the lower
bound is inflated. The tightest cheap bound is a shortest path through the
layered graph of eligible machines:
\begin{equation}
  \mathrm{LB}_j=\min_{\nu_1,\dots,\nu_{n_j}}\Bigl(\tau_{0\nu_1}+\sum_{k=1}^{n_j}p_{O_{j,k},\nu_k}
  +\sum_{k=1}^{n_j-1}\tau_{\nu_k\nu_{k+1}}\Bigr),\qquad \nu_k\in\Mset_{O_{j,k}},
\end{equation}
computed by dynamic programming in $\mathcal{O}(n_j m^2)$. It is valid with
batching, because a batch lasts at least as long as each of its members. Note
that $\sum_k\min_\nu p_{O_{j,k},\nu}+\sum_k\tau^{\min}_{O_{j,k}\to O_{j,k+1}}$ is
also valid but weaker, since it may use different machines for the same
operation in two terms.

\paragraph{References.} All CP-SAT reference makespans must be recomputed with
transport. The references of the batching experiments do not apply.

% ===========================================================================
\section{Critical assessment and open decisions}
\label{sec:critique}

\subsection{Weaknesses of the proposal}
\begin{enumerate}
  \item \textbf{Transport without resources.} Under~\ref{T:unlimited} the problem
    is an FJSP with batching and \emph{machine-dependent time lags}, not an FJSP
    with transport resources (AGVs, robots), which is the more common meaning of
    ``transport'' in the literature. This is a legitimate and cleaner variant,
    but it must be named and justified as such. The justification is that it
    isolates the spatial effect of the layout, which is what the
    representations encode, without adding a new resource type. Adding vehicles
    would require new decisions (which vehicle, empty trips) and would change
    the action space. That is outside the scope of a representation study.
  \item \textbf{Operation feature ``what was processed''.} Processed operations
    are deleted from the graph, so this feature is never observed. It is
    replaced by the incoming transport (Section~\ref{sec:repTF}).
  \item \textbf{Machine identity as a job feature.} A machine index is a label
    that only has meaning inside one instance, the same argument used to reject
    the family index. The location is given as coordinates (T-F) or as an edge
    (T-E).
  \item \textbf{Coordinate normalisation.} Without the shared scale $W$, job and
    machine coordinates live in different frames (Section~\ref{sec:repTF}). This
    is the error most likely to silently void T-F.
  \item \textbf{Machine--machine distance.} It is static, and so far removed
    from the decision that on its own it is unlikely to matter. It is kept in
    both representations (I1) for completeness, but the explicit job--machine
    transport (I3) is the main component. An ablation without I1 would show
    whether it contributes anything.
  \item \textbf{The baseline sees transport.} See Section~\ref{sec:repT0}.
    Expect smaller effects than for batching and phrase the hypothesis
    accordingly.
  \item \textbf{Batching becomes machine-dependent.} The completion rule, $\pi_j$
    and the family-node features change definition, and the batching features
    become approximations evaluated at one machine. The batch rule also stays
    myopic. It may add a partner to a batch on a machine that is far from that
    partner's next machine, increasing its future transport. The policy has no
    control over this. It is a property of the environment and must be stated.
  \item \textbf{Edge dimension.} The transport environments use width 6 instead
    of 5, so models trained for batching only cannot be reused on them. This is
    acceptable because the problem is different anyway, but the batching-only
    results and the transport results are not directly comparable.
  \item \textbf{Aggregation confound} in T-E (Section~\ref{sec:repTE}).
  \item \textbf{Calibration.} Without the $\rho$ study, a null result cannot be
    interpreted (Section~\ref{sec:grid}).
  \item \textbf{Cost.} Nine configurations with several seeds. If the budget is
    short, the most informative subset is the batching representation that won
    the batching experiments, combined with T-0, T-F and T-E.
\end{enumerate}

\subsection{Decisions to confirm}
\begin{description}[style=nextline]
  \item[D1 -- Depot.] Recommended: include a depot and the first transport. It
    gives every job a defined location at $t=0$, which the job feature needs.
    Alternative: no depot and $\tau=0$ for the first operation. This is simpler,
    but then the location of an unstarted job is undefined and needs a flag.
  \item[D2 -- Return to the depot.] Recommended: not counted. Counting it adds
    $\tau_{\mu(O_{j,n_j}),0}$ to the completion of each job and makes the choice
    of the last machine matter.
  \item[D3 -- Transport magnitude $\rho$.] Calibrated in
    Section~\ref{sec:calibration}: $\rho=0.3$.
  \item[D4 -- I3 in T-F.] Recommended: on the action edge (direct access by the
    scorer). Alternative: omit it, and let the actor derive it from the job and
    machine coordinates. This is a stricter ``node features only'' variant, but
    weaker.
  \item[D5 -- Batch features per machine.] Recommended for now: evaluate at
    $\mu^\ast_j$ / $\mu_f$ and keep the batching representation unchanged. A
    per-machine version on the \textsf{batch\_exec} edges is a later refinement.
\end{description}

% ===========================================================================
\section{Calibration of the transport magnitude}
\label{sec:calibration}

Eight batching instances of the training size (8--10 jobs, 5--10 machines, 5--6
operations per job) were solved by CP-SAT without transport and with the layouts
of four values of $\rho$. The layout is drawn after the rest of the instance, so
the jobs, families and processing times are the same for every $\rho$. All 40
solutions were proven optimal within 20~s. The \emph{increase} is
$C_{\max}(\rho)/C_{\max}(0)-1$. The \emph{share} is the fraction of the jobs'
time spent travelling in the optimal schedule,
$\sum_j(\text{transport of } j)/\sum_j C_j$.

\begin{table}[h]
\centering
\caption{Effect of $\rho$ on the optimal makespan (8 instances, mean $C_{\max}(0)=343.8$).}
\label{tab:calibration}
\begin{tabular}{c c c c}
\toprule
$\rho$ & increase (mean) & increase (min--max) & share \\
\midrule
0.1 & 0.055 & 0.029--0.102 & 0.061 \\
0.3 & 0.140 & 0.081--0.226 & 0.145 \\
0.5 & 0.198 & 0.118--0.325 & 0.170 \\
1.0 & 0.425 & 0.293--0.606 & 0.300 \\
\bottomrule
\end{tabular}
\end{table}

\noindent $\rho=0.1$ changes the makespan by about 5\%, which is within the
typical gap of a learned policy, so the representations could not be told apart.
$\rho=1$ makes transport the dominant cost (30\% of the jobs' time) and would
drown the batching signal. $\rho=0.3$ adds 14\% to the makespan, with transport
15\% of the jobs' time. It is chosen as the transport scale. The share grows
more slowly than $\rho$ (0.17 at $\rho=0.5$) because the optimal schedules
increasingly keep consecutive operations on the same or nearby machines. This is
exactly the behaviour the representations should help the policy learn. The
calibration uses only eight instances and should be repeated on the final
dataset (\texttt{python -m src.calibrate\_transport}).

% ===========================================================================
\section{Implementation}
\label{sec:impl}

\begin{description}[style=nextline]
  \item[Layout (\texttt{src/transport.py}).] Instance keys \texttt{"coords"}
    ($m\times2$ integers), \texttt{"depot"} and \texttt{"transport\_rho"}. It also
    provides the matrix $\tau$ (depot at index $m$), $\tau^{\min}$ and
    $\tau^{\mathrm{avg}}$ of~\eqref{eq:tauexp}, and the job lower bound. Instances
    without a layout have $\tau\equiv0$.
  \item[Instances (\texttt{src/batch\_generator.py}).] Option
    \texttt{transport\_rho}. Positions are uniform on $\{0,\dots,L\}^2$ with
    $L=\lceil1.5\,\rho\,\bar p\rceil$, drawn last. The dataset is built with
    \texttt{python -m src.generate\_batching\_dataset --transport-rho 0.3 --out data/transport}.
    Training instances use the $\rho$ stored in the validation dataset.
  \item[Solver and checkers.] \texttt{src/batch\_solver.py} implements
    \eqref{eq:Ta}--\eqref{eq:Td} in CP-SAT, and its \texttt{check\_schedule}
    reports transport and depot violations. \texttt{src/metrics/batching.py} adds
    the \texttt{transport} constraint to the batching checks. The job bound of
    \texttt{lower\_bound} is the shortest path of Section~\ref{sec:metrics},
    which reduces to the old bound without transport.
  \item[Environments (\texttt{src/env\_batching.py}).] The transport dynamics
    (location $\loc_j$, arrivals, transport-aware start times, batch rule with
    arrivals) are applied by every batching environment, so the three
    batching-only environments also solve transport instances correctly. A
    second class attribute \texttt{TRANSPORT\_REP} selects the representation,
    and the nine combinations are registered as
    \texttt{ojmb\_<node|edge|base>\_<t0|tf|te>}. In T-F the coordinates are the
    last two columns of the machine and job features, and \texttt{normalize\_state}
    keeps them on the shared scale, mapped to $[-1,1]$. All attributed relations
    are padded to width 6 after normalisation.
  \item[Model (\texttt{src/bopo.py}, \texttt{src/gat.py}, \texttt{src/gine.py},
    \texttt{src/transformer.py}).] The edge width of the GNN layers is taken from
    the environment (5 or 6).
  \item[Command line.] \texttt{--representation transport} trains the nine
    configurations. The test script only pairs a model with instances of its own
    kind (plain FJSP, batching, batching with transport).
\end{description}

\paragraph{Tests (\texttt{tests/test\_transport.py}).}
\begin{enumerate}[label=(\roman*)]
  \item Regression: without a layout, and with every location at the same point,
    each of the nine environments takes exactly the same schedules as the
    batching-only environment for the same actions. T-0 also has identical
    features.
  \item Random and teacher rollouts of every environment pass both checkers, and
    the makespan never falls below the lower bound.
  \item On small instances the CP-SAT solution passes the checker and is never
    worse than the teacher, and the lower bound never exceeds the solver's bound.
  \item The hand-built instance of Section~\ref{sec:example} gives the expected
    batches, arrival-based start times, T-F features and T-E edges.
  \item Violations of the transport and depot constraints are detected.
  \item A model is built for all nine representations, evaluated with the metrics
    pipeline, and trained for one BOPO update.
\end{enumerate}
The tests were checked against a mutation: when the environment ignores
transport, (ii) and (iv) fail. The existing tests (\texttt{tests/test\_metrics.py})
still pass. A two-step training run of \texttt{ojmb\_node\_tf},
\texttt{ojmb\_edge\_te} and \texttt{ojmb\_base\_t0} completed end to end.

% ===========================================================================
\section{Illustrative example}
\label{sec:example}

Three machines and a depot, with
$M_1=(0,0)$, $M_2=(4,3)$, $M_3=(10,0)$ and depot $(0,5)$:
\[
  \tau=\begin{array}{c|cccc}
     & 0 & M_1 & M_2 & M_3\\\hline
   0   & 0 & 5 & 6 & 15\\
   M_1 & 5 & 0 & 7 & 10\\
   M_2 & 6 & 7 & 0 & 9\\
   M_3 & 15& 10& 9 & 0
  \end{array}
\]
Jobs~1 and~2 have just finished: job~1 on $M_1$ at $C_1=10$ and job~2 on $M_3$
at $C_2=12$. Their current operations belong to the same family $f$, with
$B_f=2$, satisfy the order difference, and are both eligible only on $M_2$,
which is free from $F_{M_2}=15$. Arrivals: $a_{1,M_2}=10+7=17$ and
$a_{2,M_2}=12+9=21$.
\begin{itemize}
  \item Action $(M_2,1)$: $S_b=\max(15,17)=17$. Job~2 arrives at $21>17$, so it
    does not join. Job~1 runs alone, and job~2 later needs its own batch.
  \item Action $(M_2,2)$: $S_b=\max(15,21)=21$. Job~1 has arrived ($17\le21$)
    and joins. One batch, but $M_2$ is idle from 15 to 21 because of transport.
  \item Without transport, both jobs would be ready at 15 and either action would
    form the batch. Transport makes the choice of \emph{which} job opens the
    batch decisive. The batch-aware processing time $p_b$ on the action edge
    already shows the difference (a single job vs.\ the longest of the pair),
    and so do the start times ($17$ vs.\ $21$), even in T-0.
  \item After the batch, both jobs are at $M_2$. If their next operations are on
    $M_3$, both pay $\tau_{M_2M_3}=9$. T-F shows this through the incoming
    transport of those operations, and T-E through the \textsf{transfer} edges.
    T-0 shows it only once the jobs are dispatched.
\end{itemize}

\paragraph{Waiting to batch is not always worth it.} The test instance is a
variant of this example: job~0 ends on $M_1$ at 10 and arrives at $M_2$ at 17,
job~1 ends on $M_3$ at 16 and arrives at 25, and the family operations take 4
and 3. Batching them means starting at 25 and ending at 29. Running job~0 alone
(17--21) and job~1 alone (25--28) ends at 28, which equals the lower bound of
job~1 and is the CP-SAT optimum. Transport can therefore make the batch that the
completion rule would form (when job~1 opens it) the \emph{worse} choice. The
policy has to learn when to accept that.

\end{document}
